\documentclass[11pt]{article}
\usepackage{amsmath,amssymb,amsthm,mathtools}
\usepackage[T1]{fontenc}
\usepackage{lmodern}
\usepackage[margin=1in]{geometry}
\usepackage{enumitem}
\usepackage{hyperref}
\usepackage{microtype}
\hypersetup{colorlinks=true,linkcolor=blue,citecolor=blue,urlcolor=blue}

\newtheorem{theorem}{Theorem}[section]
\newtheorem{proposition}[theorem]{Proposition}
\newtheorem{lemma}[theorem]{Lemma}
\newtheorem{corollary}[theorem]{Corollary}
\newtheorem{remark}[theorem]{Remark}

\newcommand{\Z}{\mathbb Z}
\newcommand{\Sm}{S_m}
\newcommand{\Um}{U_m^{\sharp}}
\newcommand{\half}{2^{-1}}

\title{D5 121: Final-section row stitching for M5\\
The corrected-row model forces one $m^2$-cycle}
\author{}
\date{March 2026}

\begin{document}
\maketitle

\begin{abstract}
The 120 handoff reduced M5 for the color-$4$ selector surgery to the third
return
\[
U_m=(F_4^{\star})^{m^3}\big|_{P_2}
\]
on the $m^2$-state section
\[
P_2=\{(a,b,-a,0,-b):a,b\in\Z/m\Z\}.
\]
The new 120 corrected-row handoff supplies an explicit symbolic model for that
return in coordinates
\[
(a,v),\qquad v=b+c_m(a),
\]
with a sparse defect set of size $4m-7$.

This note proves the missing symbolic step: the corrected-row model itself is a
single cycle for every odd $m\ge 11$. The proof is purely finite-state and does
not use any further geometry of the torus map. The key points are:
\begin{itemize}[leftmargin=2.0em]
\item the first return of row $a=0$ is the primitive translation
      $v\mapsto v-2$;
\item the row-$0$ first-return segments are pairwise disjoint;
\item the exact return-time distribution sums to $m^2$.
\end{itemize}
Therefore the corrected-row model is one $m^2$-cycle.

As a corollary, whenever the actual return $U_m$ agrees with that symbolic
corrected-row rule, M5 holds for the color-$4$ Sel$^{\star}$ family. So the
remaining issue is no longer the final stitching theorem; it is only the
all-odd identification of the true return with the explicit corrected-row
model.
\end{abstract}

\section{The explicit corrected-row model}

Fix an odd modulus $m\ge 11$ and let
\[
\Sm=(\Z/m\Z)^2.
\]
Write states as $(a,v)$ with $a,v\in\Z/m\Z$, and let $z_t:=(0,t)$ denote the
row-$0$ state with row coordinate $0$ and column coordinate $t$.

We study the explicit model map $\Um:\Sm\to\Sm$ determined by the corrected-row
handoff. It is the clean bulk row shift
\[
(a,v)\mapsto (a+1,v)
\]
away from a sparse defect set, with the following exceptions.

\paragraph{Row $0$.}
\begin{align*}
&(0,m-8)\mapsto (0,m-10),\qquad (0,m-1)\mapsto (0,m-3),\\
&(0,v)\mapsto (1,2v+6)\qquad\text{for all other }v.
\end{align*}

\paragraph{Rows $1\le a\le 6$.}
\begin{align*}
&(1,4)\mapsto (7,2),\\
&(a,2a+2)\mapsto (a,2a),\\
&(a,2a-10)\mapsto (a,2a-12),\\
&(a,v)\mapsto (a+1,v)\qquad\text{otherwise.}
\end{align*}

\paragraph{Row $7$.}
\[
(7,16)\mapsto (7,14),
\qquad
(7,v)\mapsto (8,v)\text{ otherwise.}
\]

\paragraph{Row $8$.}
\begin{align*}
&(8,4)\mapsto (0,m-2),\\
&(8,18)\mapsto (8,16),\\
&(8,v)\mapsto (9,v)\qquad\text{otherwise.}
\end{align*}
(For $m=11,13,15,\dots$, the displayed values are always interpreted modulo
$m$.)

\paragraph{Rows $9\le a\le m-2$.}
\begin{align*}
&(a,2a+2)\mapsto (a,2a),\\
&(a,2a-12)\mapsto (a,2a-14),\\
&(a,v)\mapsto (a+1,v)\qquad\text{otherwise.}
\end{align*}

\paragraph{Row $m-1$.}
\begin{align*}
&(m-1,0)\mapsto (m-1,m-2),\\
&(m-1,m-14)\mapsto (m-1,m-16),\\
&(m-1,4)\mapsto (1,4),\\
&(m-1,2)\mapsto (0,m-1),
\end{align*}
and for all remaining row-$(m-1)$ states,
\[
(m-1,r)\mapsto
\begin{cases}
(0,(r+m-6)/2), & r\text{ odd},\\
(0,r/2-3), & r\text{ even}.
\end{cases}
\]

\begin{remark}
The exact corrected-row handoff also rewrites this as the clean bulk shift plus
five defect pieces:
\[
D_0,\quad D_{m-1},\quad L_A,\quad L_B^-,\quad L_B^+,
\]
with total size $4m-7$. The proof below only uses the row-wise form above.
\end{remark}

\section{A uniform form of the row-$(m-1)$ output}

\begin{lemma}\label{lem:last-row-unified}
For every odd $m$ and every residue
\[
r\notin\{0,m-14,2,4\},
\]
the row-$(m-1)$ rule can be written uniformly as
\[
(m-1,r)\mapsto \bigl(0,\half(r-6)\bigr)
\qquad\text{in }\Z/m\Z,
\]
where $\half=2^{-1}$ denotes the inverse of $2$ modulo $m$.
\end{lemma}

\begin{proof}
In the even branch the displayed value is
\[
r/2-3 = \half(r-6).
\]
In the odd branch the displayed value is the integer representative
\[
(r+m-6)/2,
\]
which also solves
\[
2v\equiv r-6\pmod m.
\]
Since $2$ is a unit modulo odd $m$, that solution is unique and equals
$\half(r-6)$ in $\Z/m\Z$.
\end{proof}

\section{The first return to row $0$}

Let
\[
Z_0:=\{z_t=(0,t):t\in\Z/m\Z\}
\subset \Sm
\]
be the row-$0$ section. For each $t\in\Z/m\Z$ define $\tau(t)$ to be the first
return time of $z_t$ to $Z_0$, and let
\[
\mathcal R(z_t):=(\Um)^{\tau(t)}(z_t)
\in Z_0
\]
be the first-return map.

\begin{proposition}[explicit first return]\label{prop:first-return}
For every odd $m\ge 11$ and every $t\in\Z/m\Z$,
\[
\mathcal R(z_t)=z_{t-2}.
\]
More precisely,
\[
\tau(t)=
\begin{cases}
1, & t\in\{m-8,m-1\},\\
10, & t=0,\\
2m-4, & t=1,\\
m+2, & \text{otherwise.}
\end{cases}
\]
\end{proposition}

\begin{proof}
We treat seven cases.

\smallskip
\noindent\textbf{Case 1: $t=m-8$ or $t=m-1$.}
By the row-$0$ rule,
\[
z_{m-8}\mapsto z_{m-10},
\qquad
z_{m-1}\mapsto z_{m-3}.
\]
So in either case $\tau(t)=1$ and $\mathcal R(z_t)=z_{t-2}$.

\smallskip
\noindent\textbf{Case 2: $t=0$.}
The orbit begins
\[
z_0\mapsto (1,6).
\]
No defect occurs before row $2$, where line $A$ is met:
\[
(2,6)\mapsto (2,4).
\]
After that all intermediate steps are bulk until row $8$, where the special row
$8$ rule applies:
\[
(8,4)\mapsto z_{m-2}.
\]
Thus $\tau(0)=10$ and $\mathcal R(z_0)=z_{m-2}=z_{0-2}$.

\smallskip
\noindent\textbf{Case 3: $t=1$.}
The orbit begins
\[
z_1\mapsto (1,8).
\]
The first defect is line $A$ at row $3$:
\[
(3,8)\mapsto (3,6).
\]
Then all steps are bulk until row $9$, where line $B^+$ occurs:
\[
(9,6)\mapsto (9,4).
\]
Then all steps are bulk until row $m-1$, where the special value $4$ sends the
orbit to row $1$:
\[
(m-1,4)\mapsto (1,4)\mapsto (7,2).
\]
From $(7,2)$ all steps are bulk until row $m-1$, where the special value $2$
returns to row $0$:
\[
(m-1,2)\mapsto z_{m-1}.
\]
Therefore $\tau(1)=2m-4$ and $\mathcal R(z_1)=z_{m-1}=z_{1-2}$.

\smallskip
\noindent\textbf{Case 4: $2\le t\le m-10$.}
The orbit starts at
\[
z_t\mapsto (1,\overline{2t+6}).
\]
Here $\overline{\cdot}$ denotes reduction modulo $m$.

The first defect is line $A$ at row $t+2$, because
\[
\overline{2t+6}=2(t+2)+2,
\]
and no earlier row can satisfy the same line equation. Thus
\[
(t+2,\overline{2t+6})\mapsto (t+2,\overline{2t+4}).
\]
Since $t+8\le m-2$, the next defect is line $B^+$ at row $t+8$, because
\[
\overline{2t+4}=2(t+8)-12,
\]
while all intervening rows are bulk. Hence
\[
(t+8,\overline{2t+4})\mapsto (t+8,\overline{2t+2}).
\]
After that, all remaining steps to row $m-1$ are bulk, so the orbit reaches
\[
(m-1,\overline{2t+2}).
\]
This residue is not one of the four special row-$(m-1)$ values
$0,m-14,2,4$, so Lemma~\ref{lem:last-row-unified} gives
\[
(m-1,\overline{2t+2})\mapsto z_{\half(2t+2-6)} = z_{t-2}.
\]
The step count is
\[
1+(t+1)+1+6+1+(m-t-9)+1 = m+2.
\]
Thus $\tau(t)=m+2$ and $\mathcal R(z_t)=z_{t-2}$.

\smallskip
\noindent\textbf{Case 5: $t=m-9$.}
Now
\[
z_{m-9}\mapsto (1,m-12).
\]
The first defect is line $A$ at row $m-7$:
\[
(m-7,m-12)\mapsto (m-7,m-14).
\]
All later steps to row $m-1$ are bulk, so the orbit arrives at
\[
(m-1,m-14)
\]
and then follows the special row-$(m-1)$ move
\[
(m-1,m-14)\mapsto (m-1,m-16).
\]
Applying Lemma~\ref{lem:last-row-unified} to the new residue $m-16$ gives
\[
(m-1,m-16)\mapsto z_{\half(m-22)} = z_{m-11}=z_{(m-9)-2}.
\]
The step count is
\[
1+(m-8)+1+6+1+1 = m+2.
\]
So again $\tau(t)=m+2$ and $\mathcal R(z_t)=z_{t-2}$.

\smallskip
\noindent\textbf{Case 6: $m-7\le t\le m-3$.}
The orbit starts at
\[
z_t\mapsto (1,\overline{2t+6}).
\]
Since $t\ge m-7$, the first defect is the early $B$-line at row $t+8-m\in\{1,
2,3,4,5\}$:
\[
(t+8-m,\overline{2t+6})\mapsto (t+8-m,\overline{2t+4}).
\]
After that, the next defect is line $A$ at row $t+1$:
\[
(t+1,\overline{2t+4})\mapsto (t+1,\overline{2t+2}).
\]
All later steps to row $m-1$ are bulk, so the orbit reaches
\[
(m-1,\overline{2t+2}).
\]
This is not one of the four special row-$(m-1)$ residues, hence by
Lemma~\ref{lem:last-row-unified}
\[
(m-1,\overline{2t+2})\mapsto z_{\half(2t+2-6)} = z_{t-2}.
\]
The step count is
\[
1+(t+7-m)+1+(m-7)+1+(m-t-2)+1 = m+2.
\]
Therefore $\tau(t)=m+2$ and $\mathcal R(z_t)=z_{t-2}$.

\smallskip
\noindent\textbf{Case 7: $t=m-2$.}
The orbit begins
\[
z_{m-2}\mapsto (1,2).
\]
The first defect is the early $B$-line at row $6$:
\[
(6,2)\mapsto (6,0).
\]
All later steps to row $m-1$ are bulk, so the orbit reaches $(m-1,0)$ and then
follows the special row-$(m-1)$ move
\[
(m-1,0)\mapsto (m-1,m-2).
\]
Applying Lemma~\ref{lem:last-row-unified} to the residue $m-2$ gives
\[
(m-1,m-2)\mapsto z_{\half(m-8)} = z_{m-4}=z_{(m-2)-2}.
\]
The step count is
\[
1+5+1+(m-7)+1+1 = m+2.
\]
So $\tau(m-2)=m+2$ and $\mathcal R(z_{m-2})=z_{m-4}$.

These seven cases exhaust all residues $t\in\Z/m\Z$.
\end{proof}

\section{Disjoint first-return segments}

For each $t\in\Z/m\Z$, define the first-return segment
\[
S_t:=\{(\Um)^j(z_t):0\le j<\tau(t)\}.
\]
Thus $S_t$ contains the starting row-$0$ state $z_t$, contains no other
row-$0$ states, and ends just before the next row-$0$ hit
$\mathcal R(z_t)=z_{t-2}$.

\begin{lemma}\label{lem:segments-disjoint}
The family $(S_t)_{t\in\Z/m\Z}$ is pairwise disjoint.
\end{lemma}

\begin{proof}
Suppose $x\in S_t\cap S_u$.

If $x\in Z_0$, then by construction $x$ must be the initial row-$0$ point of
both segments. Hence $x=z_t=z_u$, so $t=u$.

Now assume $x\notin Z_0$. From $x$, the future orbit until the next hit of
$Z_0$ is uniquely determined. Since $x\in S_t$, that next row-$0$ hit is
$\mathcal R(z_t)=z_{t-2}$. Since $x\in S_u$, the same next row-$0$ hit is also
$\mathcal R(z_u)=z_{u-2}$. Therefore $z_{t-2}=z_{u-2}$, hence $t=u$.

So $S_t\cap S_u\ne\varnothing$ implies $t=u$.
\end{proof}

\section{The corrected-row model is one cycle}

\begin{theorem}[row stitching theorem]\label{thm:model-one-cycle}
For every odd $m\ge 11$, the corrected-row model $\Um$ is one cycle on the full
$m^2$-state set $\Sm$.
\end{theorem}

\begin{proof}
By Proposition~\ref{prop:first-return}, the first-return map on row $0$ is the
translation
\[
\mathcal R(z_t)=z_{t-2}.
\]
Since $-2$ is a unit modulo odd $m$, $\mathcal R$ is one $m$-cycle on the row
$0$ section $Z_0$.

Now compute the total length of the first-return segments. There are exactly two
residues with return time $1$, namely $m-8$ and $m-1$; one residue with return
time $10$, namely $0$; one residue with return time $2m-4$, namely $1$; and the
remaining $m-4$ residues have return time $m+2$. Therefore
\[
\sum_{t\in\Z/m\Z}\lvert S_t\rvert
=2\cdot 1 + 10 + (2m-4) + (m-4)(m+2) = m^2.
\]
By Lemma~\ref{lem:segments-disjoint}, the segments $S_t$ are pairwise disjoint.
Their union therefore has cardinality $m^2$, which is the full size of $\Sm$.

Finally, because the row-$0$ return is the single cycle
\[
z_0\to z_{-2}\to z_{-4}\to\cdots\to z_{-2(m-1)}\to z_0,
\]
concatenating the segments in that order produces one closed orbit of length
$m^2$. Since that orbit already covers all of $\Sm$, the map $\Um$ is a single
$m^2$-cycle.
\end{proof}

\section{Consequence for the actual M5 target}

Let $U_m=(F_4^{\star})^{m^3}|_{P_2}$ be the actual third return from the 120
M5 setup, and let $v=b+c_m(a)$ be the corrected-row coordinate from the 120
handoff.

\begin{corollary}[conditional M5 closure from the corrected-row model]
Assume that, for a given odd $m\ge 11$, the actual return $U_m$ agrees in the
corrected-row coordinates $(a,v)$ with the symbolic model $\Um$ from
Section~1. Then $U_m$ is one cycle on $P_2$.

Consequently, by the 120 nested-return theorem, the color-$4$ torus map
$F_4^{\star}$ is Hamilton.
\end{corollary}

\begin{proof}
Under the stated hypothesis, the corrected-row coordinate change identifies the
actual return $U_m$ with the explicit model $\Um$. By
Theorem~\ref{thm:model-one-cycle}, that model is one cycle. So $U_m$ is one
cycle on $P_2$. The final statement is exactly the reduction theorem proved in
120.
\end{proof}

\begin{remark}
This note proves the final stitching theorem once the corrected-row model is in
hand. So the remaining all-odd issue for full M5 is no longer the final section
cycle structure. It is only the all-odd derivation of the exact corrected-row
rule for the true return $U_m$.
\end{remark}

\end{document}
